\item \textbf{¿Por qué no es posible la siguiente situación?}
Dos rendijas estrechas están separadas por $8.00\text{ mm}$ en una pieza de metal. Un haz de microondas golpea el metal perpendicularmente, pasa a través de las dos rendijas y luego procede hacia una pared a cierta distancia. La longitud de onda de la radiación se estima en $1.00\text{ cm} \pm 5\%$. Al mover un detector de microondas a lo largo de la pared para estudiar el patrón de interferencia, se logra medir con alta precisión la posición de la franja brillante $m = 1$, obteniendo una medición exitosa de la longitud de onda de la radiación.